% !TeX root = ./main.tex

\chapter*{\large ЗАКЛЮЧЕНИЕ}
\addcontentsline{toc}{chapter}{ЗАКЛЮЧЕНИЕ}

\textbf{Основные научные и практические результаты работы}

\begin{enumerate}
    \item Систематизированы современные информационные технологии и теоретические подходы к переносу графовых динамических моделей между дискретизациями (спектральные фильтры и трансформеры, предельные графонные ОДУ, комбинаторные сети реакции-адвекции-диффузии и континуальные нейронные операторы). Показано, что если на молекулярных графах фиксированной ковалентной топологии одношаговый RBF-вентиль гомогенности эффективен (архитектура GRADE~[\ref{itm:a_grade}]), то при измельчении пространственных сеток изображений ($h_G\to 0$) в комбинаторных моделях типа AdaCDE он вырождается по закону $\Theta(h_G^2)\to 0$, блокируя направленный перенос на высоком разрешении~[\ref{itm:a1}--\ref{itm:a2}].
    \item Разработан и программно реализован на открытом стеке Python~3.11 (\texttt{NumPy}, \texttt{SciPy CSR/CSC}) континуально-согласованный класс разреженных операторов \texttt{CDQuad}, объединяющий кусочно-линейную конечно-элементную жёсткость анизотропного тензора диффузии $K(x)$, строго консервативные противопоточные потоки на барицентрических дуальных ячейках Дональда и нелокальный гауссов вентиль фиксированного физического радиуса $r_0>0$. Установлено, что бездивергентный перенос не входит в показатель экспоненты Гронуолла, обеспечивая равномерную устойчивость траекторий на длинных горизонтах времени и строго линейный масштаб ошибки по числу Пекле ($R^2=0{,}9998$ при $\Pe\in[0,100]$)~[\ref{itm:a1}].
    \item Предложен быстровычислимый ($4{,}2$~мс при $N=3025$) метод сертификации ошибки дискретизации без запуска эталонного решателя, основанный на эллиптическом двойственном словарном дефекте $\widehat{\delta}_{-1}(G)$ в норме $H_G^{-1}$ и раздельной конформной калибровке. Показано, что переход от поточечной $L^2$-невязки к двойственной норме повышает корреляцию с истинной ошибкой переноса с $r=0{,}5590$ до $r=0{,}9803$, сужает ширину $90\,\%$-го доверительного интервала на $36{,}1\,\%$ и сохраняет $100\,\%$-е покрытие при глобальном анизотропном сдвиге сеток, а глобально липшицево сглаживание порога позволяет переносить гарантии покрытия непосредственно на площадь и барицентр структуры.
    \item На синтетическом бенчмарке мигрирующих структур и реальной цейтраферной флуоресцентной микроскопии ядерного транспорта живой клетки (\texttt{NPCsingleNucleus.tif}) подтверждена эффективность регламента «обучение на грубом биннинге~-- измерение без дообучения на высоком разрешении»: при росте числа узлов в $8{,}4\text{--}9$~раз ошибка поля \texttt{CDQuad} снижается в $2{,}28$~раза (превосходя AdaCDE в $29{,}1$~раза по полю и в $45{,}7$~раза по скорости фронта), а межмасштабный дрейф на живой микроскопии сокращается в $4{,}1\text{--}4{,}4$~раза. Чётко локализованы объективные границы применимости метода на несимплициальных графах суперпикселей SLIC и при локальных «слепых пятнах» датчика.
\end{enumerate}

\vspace{0.35cm}
\textbf{Основные тенденции развития исследуемой области}

Обобщение результатов обзора и вычислительных экспериментов выявляет три ведущие тенденции развития информационных технологий пространственно-временного моделирования: 1)~переход от эмпирических комбинаторных ГНС к гибридным континуально-согласованным операторам, явно использующим метрические квадратуры ячеек и граней; 2)~интеграция методов непараметрической конформной калибровки непосредственно в конвейеры биомедицинских и инженерных измерений для выдачи сертифицированных доверительных интервалов в реальном времени; 3)~ориентация на воспроизводимые легковесные открытые программные комплексы, способные функционировать локально на стандартных рабочих станциях без облачных GPU-ресурсов.

\vspace{0.35cm}
\textbf{Рекомендации по практическому использованию результатов}

Открытый программный комплекс \texttt{CDQuad} и методика статистической сертификации морфометрических дескрипторов рекомендуются к внедрению в профильных лабораториях биомедицинской информатики и оптической микроскопии (для щадящей калибровки динамических протоколов съёмки живых клеток на режимах аппаратного биннинга с последующим анализом кадров полного разрешения), в системах экологического и инженерного мониторинга на нерегулярных сенсорных триангуляциях, а также в учебном процессе факультета прикладной математики и информатики БГУ при преподавании дисциплин по компьютерному анализу изображений, статистическому машинному обучению и численным методам математической физики.

\newpage
